% No numbered results (statements-of-record rows resolved here: none).
\section{Discussion}
\label{sec:discussion}

\paragraph{What the route requires.}
The duality-confinement route reduces the Riemann hypothesis to upper
bounds tending to zero for the window sums \(A_W\), or for operators
dominating the associated ledgers.  For windows closed under the
reflection \(s\mapsto1-\bar s\) (\Cref{thm:rh:concrete-window-dc}), or under the map \(s\mapsto1-s\) with invariant weights
(\Cref{prop:rh:window-equivalences}), and along windows that eventually
contain every zero, the existence of such bounds is the Riemann
hypothesis itself, whatever form the bounds take among those
considered.  What would be new is a source of bounds: an argument, for
instance from an explicit formula or a trace formula, that produces
them from zeta-function data without presupposing the conclusion.
\Cref{prop:rh:window-obstructions} records the elementary fact that an operator-level transport bound for
the identity target, of the form stated there, cannot tend to zero on a
nonempty window, so a payment of that kind must be evaluated on the
actual displacement vector.

\paragraph{Relation to classical programs.}
Several classical approaches seek a positivity statement equivalent to
the Riemann hypothesis: Weil's criterion, which asks for positivity of
a quadratic functional built from the explicit formula
\citep{Weil1952,Bombieri2000}, and Li's criterion \citep{Li1997}; the Nyman--Beurling criterion
\citep{Nyman1950,Beurling1955,BaezDuarte2003}; de Branges' spaces of
entire functions \citep{deBranges1968,deBranges1986}; and spectral
interpretations of the zeros \citep{BerryKeating1990,Connes1999}.  The window sums of this paper are also positivity data, but of a much
cruder kind: they are the squared displacements themselves, and, unlike Weil's functional, which is built from the explicit
formula relating zeros to primes, they involve no primes.  A bound
for them that did not presuppose the Riemann hypothesis would presumably
have to come from a deeper source of this kind; nothing in the present paper
supplies one.

\paragraph{Relation to the companion papers.}
The duality-confinement theorem is proved in \citet{TsiokosSDTC2026};
this paper is its application to the zeros of \(\zeta\), and
\Cref{thm:rh:concrete-window-dc} is the zeta-specific case of the
converse proved there.  The residual operator of \Cref{sec:finite_windows} is the
identity-audit case of the adequacy residual of
\citet{TsiokosXiCompanion2026}, which in the companion Navier--Stokes
paper \citep{TsiokosNavier2026} is applied to a physical target; the
metatheory paper \citep{TsiokosMetaTheory2026} relates the two
applications through a common Gaussian observation.  None of these
connections supplies a payment for either problem.

\paragraph{Further directions.}
Three questions remain open: to construct a recognition source from a
trace formula, which would turn \Cref{thm:rh:full-classical} into a
proof; to verify the framework's admissibility conditions for such a
construction; and to treat other \(L\)-functions, each of which needs
its own ledger and source.
